% =============================================================================
% Conclusion générale
% =============================================================================
\btpLiminary{Conclusion générale}

\section*{Résumé des contributions}

Ce travail de recherche a contribué sur cinq plans.

\begin{enumerate}
  \item \btpterm{Analyse du domaine.} Une caractérisation du processus de
        réponse aux appels d'offres dans le BTP camerounais, de sa
        production documentaire et de ses points de friction~: relecture
        répétée, reconstitution du patrimoine, contrôle de complétude tardif.
  \item \btpterm{État de l'art.} Un panorama raisonné des disciplines
        mobilisées --- traitement automatique du langage, recherche
        d'information, génération augmentée par récupération, compréhension
        de documents, résolution d'entités, workflows et supervision
        humaine --- en distinguant les concepts des outils.
  \item \btpterm{Méthodologie.} Une approche structurée, fondée sur deux
        horloges, une recherche hybride, un découpage parent-enfant, une
        génération par section et une revue humaine constitutive.
  \item \btpterm{Architecture.} Une plateforme modulaire qui incarne ces
        choix, évaluable et substituable pièce par pièce.
  \item \btpterm{Validation.} Une première campagne empirique qui soutient
        les hypothèses H1 à H5 et démontre la faisabilité opérationnelle
        du système.
\end{enumerate}

Le résultat n'est pas « un modèle qui rédige des offres ». C'est une
chaîne de traitement dans laquelle l'intelligence artificielle accélère
la recherche et la rédaction, et dans laquelle l'ingénieur conserve la
responsabilité de l'acte engageant.

\section*{Limites du travail}

Les limites énoncées au chapitre~\ref{chap:resultats} valent pour
l'ensemble du mémoire. La campagne d'évaluation reste de taille
modeste~; le graphe de connaissances n'est pas activé~; le français
juridique camerounais demeure un terrain difficile pour des modèles
entraînés ailleurs. Ces limites bornent la portée des conclusions, elles
ne les invalident pas.

\section*{Perspectives}

Cinq perspectives se dessinent.

\begin{enumerate}
  \item \btpterm{Industrialisation.} Passage à une solution multi-organisations,
        avec cloisonnement, journalisation d'audit et engagements de
        service.
  \item \btpterm{Graphe de connaissances.} Activation de la phase 5, une
        fois le RAG hybride consolidé, pour les questions relationnelles.
  \item \btpterm{Intelligence localement agentique.} Encapsuler une
        recherche multi-sauts dans un n\oe ud du workflow, sans rendre
        l'ensemble autonome.
  \item \btpterm{Multimodalité.} Intégration des plans, schémas et images,
        aujourd'hui absents de l'index.
  \item \btpterm{Adaptation de domaine.} Affinage ou spécialisation des
        modèles au français des marchés publics et au vocabulaire du BTP.
\end{enumerate}

Au-delà de ces prolongements techniques, le bénéfice attendu est
organisationnel~: transformer un patrimoine documentaire dormant en une
mémoire d'entreprise effectivement mobilisable, dossier après dossier,
sous le contrôle de ceux qui en portent la responsabilité.

\btpEndMark
\clearpage
